% ╔══════════════════════════════════════════════════╗
% ║  CV – İlan Eşleşmesi: Teknik Rapor              ║
% ║  Dil: Türkçe  |  Derleyici: pdflatex (Overleaf) ║
% ╚══════════════════════════════════════════════════╝

\documentclass[12pt, a4paper]{article}

% ── Temel Paketler ───────────────────────────────────────────
\usepackage[utf8]{inputenc}
\usepackage[T1]{fontenc}
\usepackage[turkish]{babel}
\usepackage{geometry}
\usepackage{amsmath}
\usepackage{amssymb}
\usepackage{graphicx}
\usepackage{booktabs}
\usepackage{array}
\usepackage{xcolor}
\usepackage{hyperref}
\usepackage{listings}
\usepackage{fancyhdr}
\usepackage{titlesec}
\usepackage{enumitem}
\usepackage{float}
\usepackage{caption}
\usepackage{mdframed}
\usepackage{lmodern}
\usepackage{microtype}

% ── Sayfa Düzeni ─────────────────────────────────────────────
\geometry{
    top=2.5cm, bottom=2.5cm,
    left=2.5cm, right=2.5cm,
    headheight=15pt
}

% ── Renkler ──────────────────────────────────────────────────
\definecolor{anarenk}{RGB}{233,69,96}
\definecolor{morrenk}{RGB}{99,102,241}
\definecolor{yesilenk}{RGB}{74,222,128}
\definecolor{grirenk}{RGB}{160,174,192}
\definecolor{kodbg}{RGB}{26,31,46}
\definecolor{turuncu}{RGB}{245,158,11}

% ── Bağlantılar ──────────────────────────────────────────────
\hypersetup{
    colorlinks=true,
    linkcolor=anarenk,
    urlcolor=morrenk,
    citecolor=morrenk,
    pdftitle={CV – İlan Eşleşmesi: Teknik Rapor},
    pdfsubject={İnsan Kaynakları, NLP, Makine Öğrenmesi}
}

% ── Kod Stili ────────────────────────────────────────────────
\lstdefinestyle{kod}{
    basicstyle=\ttfamily\small,
    keywordstyle=\color{morrenk}\bfseries,
    stringstyle=\color{yesilenk},
    commentstyle=\color{grirenk}\itshape,
    numberstyle=\tiny\color{grirenk},
    numbers=left,
    numbersep=8pt,
    backgroundcolor=\color{kodbg},
    frame=single,
    framerule=0.4pt,
    rulecolor=\color{grirenk},
    breaklines=true,
    tabsize=4,
    showstringspaces=false,
    xleftmargin=12pt,
    basicstyle=\ttfamily\small\color{white},
}

% ── Bilgi Kutusu (mdframed ile — güvenli) ────────────────────
\mdfdefinestyle{bilgikutsu}{
    linecolor=anarenk,
    linewidth=1.5pt,
    backgroundcolor=kodbg,
    fontcolor=white,
    innertopmargin=8pt,
    innerbottommargin=8pt,
    innerleftmargin=10pt,
    innerrightmargin=10pt,
    roundcorner=4pt
}

% ── Başlık Stilleri ──────────────────────────────────────────
\titleformat{\section}
    {\Large\bfseries\color{anarenk}}
    {\thesection}{1em}{}[\color{anarenk}\titlerule]

\titleformat{\subsection}
    {\large\bfseries\color{morrenk}}
    {\thesubsection}{1em}{}

\titleformat{\subsubsection}
    {\normalsize\bfseries}
    {\thesubsubsection}{1em}{}

% ── Üstbilgi / Altbilgi ──────────────────────────────────────
\pagestyle{fancy}
\fancyhf{}
\fancyhead[L]{\small\color{grirenk}CV -- İlan Eşleşmesi}
\fancyhead[R]{\small\color{grirenk}İnsan Kaynakları $\cdot$ NLP}
\fancyfoot[C]{\small\color{grirenk}\thepage}
\renewcommand{\headrulewidth}{0.4pt}

% ─────────────────────────────────────────────────────────────
\begin{document}
% ─────────────────────────────────────────────────────────────

% ── Kapak Sayfası ────────────────────────────────────────────
\begin{titlepage}
    \centering
    \vspace*{2cm}

    {\Huge\bfseries\color{anarenk} CV -- İlan Eşleşmesi}\\[0.5cm]
    {\Large\color{morrenk} İnsan Kaynakları Projesi: Teknik Rapor}\\[2cm]

    \begin{mdframed}[style=bilgikutsu]
    \centering
    \large\color{white}
    \textbf{Proje Özeti}\\[0.4cm]
    \normalsize\color{grirenk}
    Bu rapor, aday CV'lerinin iş ilanlarıyla metinsel uyumunu\\
    ölçen bir NLP sisteminin teknik belgesidir.\\[0.4cm]
    \color{white}
    \textbf{Yöntem:} TF-IDF Vektörleştirme + Kosinüs Benzerliği\\
    \textbf{Arayüz:} Streamlit Dashboard\\
    \textbf{Dil:} Python 3.11+
    \end{mdframed}

    \vspace{2cm}

    \begin{tabular}{ll}
        \toprule
        \textbf{Proje Adı}     & CV -- İlan Eşleşmesi \\
        \textbf{Alan}          & İnsan Kaynakları / NLP \\
        \textbf{Teknolojiler}  & Python, Scikit-learn, Streamlit, Plotly \\
        \textbf{Sürüm}         & v1.0 \\
        \textbf{Tarih}         & \today \\
        \bottomrule
    \end{tabular}

    \vfill
    {\small\color{grirenk} Bu rapor proje gelişimi boyunca güncellenecektir.}
\end{titlepage}

% ── İçindekiler ──────────────────────────────────────────────
\tableofcontents
\newpage

% ─────────────────────────────────────────────────────────────
\section{Giriş ve Motivasyon}
% ─────────────────────────────────────────────────────────────

\subsection{Problem Tanımı}

Günümüz iş dünyasında insan kaynakları departmanları, her iş ilanı için onlarca hatta yüzlerce CV değerlendirmek zorunda kalmaktadır. Bu süreç zaman alıcı, maliyetli ve tutarsız sonuçlar doğurabilmektedir. \textbf{CV -- İlan Eşleşmesi} projesi, bu sorunu doğal dil işleme (NLP) teknikleriyle çözmeyi amaçlamaktadır.

\begin{mdframed}[style=bilgikutsu]
\color{white}
\textbf{Temel Problem:} Aday CV'lerinin iş ilanlarıyla \textbf{metinsel uyumunu otomatik olarak ölçmek} ve en uygun adayları sıralamak.
\end{mdframed}

\subsection{Projenin Kapsamı}

Bu proje şu bileşenleri kapsamaktadır:

\begin{itemize}
    \item \textbf{Metin Ön İşleme:} Unicode normalizasyonu, stopword eleme, noktalama temizleme
    \item \textbf{Vektörleştirme:} TF-IDF (birincil) ve Sentence-Transformers (opsiyonel)
    \item \textbf{Benzerlik Hesaplama:} Kosinüs Benzerliği formülü
    \item \textbf{Sıralama:} Top-N aday çıktısı
    \item \textbf{Görselleştirme:} Streamlit tabanlı interaktif dashboard
\end{itemize}

% ─────────────────────────────────────────────────────────────
\section{Matematiksel Altyapı}
% ─────────────────────────────────────────────────────────────

\subsection{TF-IDF (Term Frequency -- Inverse Document Frequency)}

TF-IDF, bir kelimenin belgede ne kadar önemli olduğunu ölçen istatistiksel bir yöntemdir.

\subsubsection{Term Frequency (TF)}

\begin{equation}
    \text{TF}(t,\, d) = \frac{f_{t,d}}{\displaystyle\sum_{t' \in d} f_{t',d}}
\end{equation}

Burada $f_{t,d}$, $t$ teriminin $d$ belgesinde geçme sayısını ifade eder.

\subsubsection{Inverse Document Frequency (IDF)}

\begin{equation}
    \text{IDF}(t,\, D) = \log \frac{|D|}{\left|\{d \in D : t \in d\}\right|}
\end{equation}

\subsubsection{TF-IDF Ağırlığı}

\begin{equation}
    \text{TF-IDF}(t,\, d,\, D) \;=\; \text{TF}(t,\, d) \;\times\; \text{IDF}(t,\, D)
\end{equation}

Bu projede \texttt{sublinear\_tf=True} parametresi kullanılarak log-ölçekli TF uygulanmaktadır:

\begin{equation}
    \text{TF}_{\log}(t,\, d) = 1 + \log(f_{t,d})
\end{equation}

\subsection{Kosinüs Benzerliği}

İki metin vektörü $\mathbf{A}$ ve $\mathbf{B}$ arasındaki kosinüs benzerliği:

\begin{equation}
    \boxed{
        \text{benzerlik}(\mathbf{A},\, \mathbf{B})
        \;=\;
        \cos(\theta)
        \;=\;
        \frac{\mathbf{A} \cdot \mathbf{B}}{\|\mathbf{A}\| \cdot \|\mathbf{B}\|}
        \;=\;
        \frac{\displaystyle\sum_{i=1}^{n} A_i B_i}
             {\displaystyle\sqrt{\sum_{i=1}^{n} A_i^2}
              \;\cdot\;
              \sqrt{\sum_{i=1}^{n} B_i^2}}
    }
\end{equation}

\begin{itemize}
    \item Sonuç aralığı: $[0,\; 1]$
    \item $\cos(\theta) \to 1$ \;$\Rightarrow$\; tam uyum
    \item $\cos(\theta) \to 0$ \;$\Rightarrow$\; uyumsuzluk
\end{itemize}

\subsection{Skor Yorumlama Tablosu}

\begin{table}[H]
\centering
\caption{Kosinüs Benzerlik Skoru Yorumlama Rehberi}
\begin{tabular}{ccc}
\toprule
\textbf{Skor Aralığı} & \textbf{Değerlendirme} & \textbf{Öneri} \\
\midrule
$0{,}60 \leq s \leq 1{,}00$ & {\color{yesilenk}Yüksek Uyum}  & Mülakat çağrısı \\
$0{,}35 \leq s < 0{,}60$    & {\color{turuncu}Orta Uyum}     & Detaylı inceleme \\
$0{,}00 \leq s < 0{,}35$    & {\color{red}Düşük Uyum}        & Elenebilir \\
\bottomrule
\end{tabular}
\end{table}

% ─────────────────────────────────────────────────────────────
\section{Sistem Mimarisi}
% ─────────────────────────────────────────────────────────────

\subsection{Dosya Yapısı}

\begin{lstlisting}[style=kod, language={}, caption={Proje Klasör Yapısı}]
CV_MATCHING/
|-- app.py                           # Streamlit ana uygulamasi
|-- matching_engine.py               # NLP motoru
|-- sample_data.py                   # Ornek veriler
|-- test_matching.py                 # Terminal testi
|-- requirements.txt                 # Bagimliliklar
|-- Dockerfile                       # Docker imaji
|-- docker-compose.yml               # Docker Compose
|-- notebooks/
|   |-- 01_eda_ve_onisleme.ipynb     # Keşifsel analiz
|   |-- 02_model_karsilastirma.ipynb # Model karsilastirmasi
|-- rapor/
|   |-- cv_ilan_eslesme_rapor.tex    # Bu rapor (LaTeX)
|-- cvs_ornek/
    |-- aday1_cv.txt
    |-- aday2_cv.txt
\end{lstlisting}

\subsection{Bileşen Açıklamaları}

\begin{table}[H]
\centering
\caption{Proje Dosyaları ve Görevleri}
\begin{tabular}{lp{9cm}}
\toprule
\textbf{Dosya} & \textbf{Açıklama} \\
\midrule
\texttt{app.py}             & Streamlit arayüzü -- tarayıcıda açılan dashboard \\
\texttt{matching\_engine.py}& Metin ön işleme, TF-IDF, Cosine Similarity mantığı \\
\texttt{sample\_data.py}    & 5 örnek CV ve 1 örnek iş ilanı \\
\texttt{test\_matching.py}  & Streamlit olmadan terminalde hızlı test \\
\texttt{requirements.txt}   & Python bağımlılık listesi \\
\texttt{Dockerfile}         & Docker konteyner imajı \\
\texttt{notebooks/01\_*}    & EDA ve ön işleme adımları \\
\texttt{notebooks/02\_*}    & TF-IDF -- Sentence-Transformers karşılaştırması \\
\bottomrule
\end{tabular}
\end{table}

\subsection{Metin İşleme Pipeline'ı}

Ham metin aşağıdaki 6 adımdan geçirilir:

\[
\text{Ham Metin}
\xrightarrow{(1)}
\text{Unicode}
\xrightarrow{(2)}
\text{Küçük Harf}
\xrightarrow{(3)}
\text{URL/Email}
\xrightarrow{(4)}
\text{Sayılar}
\xrightarrow{(5)}
\text{Noktalama}
\xrightarrow{(6)}
\text{Stopword}
\;\longrightarrow\;
\text{Token}
\]

% ─────────────────────────────────────────────────────────────
\section{Teknik Uygulama}
% ─────────────────────────────────────────────────────────────

\subsection{Bağımlılıklar}

\begin{table}[H]
\centering
\caption{Kullanılan Python Kütüphaneleri}
\begin{tabular}{lll}
\toprule
\textbf{Kütüphane} & \textbf{Versiyon} & \textbf{Kullanım Amacı} \\
\midrule
streamlit            & $\geq$ 1.35.0 & Web arayüzü \\
scikit-learn         & $\geq$ 1.4.0  & TF-IDF ve Cosine Similarity \\
pandas               & $\geq$ 2.0.0  & Veri işleme \\
plotly               & $\geq$ 5.18.0 & İnteraktif grafikler \\
nltk                 & $\geq$ 3.8.0  & Stopword listesi \\
numpy                & $\geq$ 1.26.0 & Vektör işlemleri \\
\midrule
\multicolumn{3}{c}{\textit{Opsiyonel}} \\
\midrule
sentence-transformers & $\geq$ 2.7.0 & Derin anlam vektörleri \\
torch                 & $\geq$ 2.2.0 & Sentence-Transformers için \\
\bottomrule
\end{tabular}
\end{table}

\subsection{Temel Kod -- Kosinüs Benzerliği}

\begin{lstlisting}[style=kod, language=Python,
    caption={Kosinüs Benzerliği -- Sıfırdan Uygulama}]
import numpy as np

def cosine_similarity_manuel(A: np.ndarray,
                              B: np.ndarray) -> float:
    """
    Formul: cos(theta) = (A . B) / (||A|| * ||B||)
    """
    nokta_carpim = np.dot(A, B)       # A . B
    norm_A = np.linalg.norm(A)        # ||A||
    norm_B = np.linalg.norm(B)        # ||B||

    if norm_A == 0 or norm_B == 0:
        return 0.0   # Sifir vektor korumasi

    return float(nokta_carpim / (norm_A * norm_B))
\end{lstlisting}

\subsection{TF-IDF Vektörleştirici}

\begin{lstlisting}[style=kod, language=Python,
    caption={TF-IDF Vektörleştirici -- Scikit-learn}]
from sklearn.feature_extraction.text import TfidfVectorizer
from sklearn.metrics.pairwise import cosine_similarity

vec = TfidfVectorizer(
    ngram_range=(1, 2),   # unigram + bigram
    max_features=5000,
    sublinear_tf=True     # log olcekli TF
)

tum_metinler = [is_ilani_metni] + cv_metinleri
matrix = vec.fit_transform(tum_metinler).toarray()

is_vec = matrix[0]
cv_vecs = matrix[1:]
skorlar = cosine_similarity([is_vec], cv_vecs)[0]
\end{lstlisting}

% ─────────────────────────────────────────────────────────────
\section{Kurulum ve Çalıştırma}
% ─────────────────────────────────────────────────────────────

\subsection{Gereksinimler}

\begin{itemize}
    \item Python 3.9 veya üzeri
    \item pip (Python paket yöneticisi)
    \item İsteğe bağlı: Docker Desktop
\end{itemize}

\subsection{Kurulum Adımları}

\begin{lstlisting}[style=kod, language={}, caption={Terminal Komutları}]
# 1. Klasor olustur
mkdir cv_matching && cd cv_matching

# 2. Sanal ortam olustur
python -m venv venv
source venv/bin/activate   # Linux/Mac
# venv\Scripts\activate    # Windows

# 3. Bagimliliklar
pip install -r requirements.txt

# 4. Test (terminal)
python test_matching.py

# 5. Streamlit Dashboard
streamlit run app.py
# Erisim: http://localhost:8501
\end{lstlisting}

\subsection{Docker ile Çalıştırma}

\begin{lstlisting}[style=kod, language={}, caption={Docker Komutları}]
# Imaj olustur
docker build -t cv-matching .

# Calistir
docker run -p 8501:8501 cv-matching

# Ya da Docker Compose ile
docker compose up
\end{lstlisting}

% ─────────────────────────────────────────────────────────────
\section{Sonuçlar ve Değerlendirme}
% ─────────────────────────────────────────────────────────────

\subsection{Örnek Çıktı}

E-ticaret sektörüne yönelik ``Senior Data Scientist'' ilanı için 5 aday arasındaki örnek sıralama:

\begin{table}[H]
\centering
\caption{Örnek Aday Sıralama Sonuçları}
\begin{tabular}{clcc}
\toprule
\textbf{Sıra} & \textbf{Aday Adı} & \textbf{Skor (\%)} & \textbf{Değerlendirme} \\
\midrule
1 & Laura Martínez  & 72.4 & {\color{yesilenk}Yüksek} \\
2 & Carlos Romero   & 58.1 & {\color{turuncu}Orta} \\
3 & Daniela Torres  & 41.3 & {\color{turuncu}Orta} \\
4 & Ivan Perez      & 24.1 & {\color{red}Düşük} \\
5 & Sofia Vargas    & 11.8 & {\color{red}Düşük} \\
\bottomrule
\end{tabular}
\end{table}

\subsection{Güçlü ve Zayıf Yönler}

\begin{table}[H]
\centering
\caption{Sistem Değerlendirmesi}
\begin{tabular}{p{6.5cm}p{6.5cm}}
\toprule
\textbf{Güçlü Yönler} & \textbf{Geliştirme Alanları} \\
\midrule
Hızlı sonuç (TF-IDF $<$ 0.1 sn) & Anlam körü (eş anlamlı kelimeler) \\
Kolay yorumlanabilir              & PDF desteği yok (sadece .txt) \\
Türkçe + İngilizce destek        & ATS entegrasyonu eksik \\
Streamlit ile kolay kullanım     & Gerçek veri ile test edilmedi \\
\bottomrule
\end{tabular}
\end{table}

% ─────────────────────────────────────────────────────────────
\section{Gelecek Çalışmalar}
% ─────────────────────────────────────────────────────────────

\begin{enumerate}
    \item \textbf{PDF Desteği:} \texttt{pdfplumber} ile CV PDF dosyalarını otomatik okuma
    \item \textbf{Derin Öğrenme:} Sentence-Transformers (\texttt{all-MiniLM-L6-v2}) modeli varsayılan yapma
    \item \textbf{ATS Entegrasyonu:} Greenhouse, Workable, BambooHR API bağlantısı
    \item \textbf{Docker Dağıtımı:} Üretim ortamında konteyner tabanlı çalışma
    \item \textbf{Streamlit Cloud:} Halka açık URL ile ücretsiz erişim
    \item \textbf{Çoklu İlan:} Aynı anda birden fazla iş ilanı ile karşılaştırma
    \item \textbf{Açıklanabilirlik:} SHAP ile hangi kelimelerin skoru etkilediğini gösterme
\end{enumerate}

% ─────────────────────────────────────────────────────────────
\section{Değişiklik Günlüğü}
% ─────────────────────────────────────────────────────────────

\begin{table}[H]
\centering
\caption{Proje Sürüm Geçmişi}
\begin{tabular}{llp{7.5cm}}
\toprule
\textbf{Sürüm} & \textbf{Tarih} & \textbf{Değişiklikler} \\
\midrule
v1.0 & \today
     & İlk sürüm: TF-IDF + Cosine Similarity, Streamlit UI, 5 örnek aday,
       metin ön işleme pipeline'ı, .txt dosya yükleme, CSV export,
       2 Jupyter notebook, LaTeX raporu \\
\midrule
v1.1 & -- & Planlanan: PDF desteği, çoklu ilan karşılaştırması \\
v1.2 & -- & Planlanan: Sentence-Transformers varsayılan mod \\
v2.0 & -- & Planlanan: ATS API entegrasyonu, Docker dağıtımı \\
\bottomrule
\end{tabular}
\end{table}

% ─────────────────────────────────────────────────────────────
\section*{Ekler}
\addcontentsline{toc}{section}{Ekler}
% ─────────────────────────────────────────────────────────────

\subsection*{A. Tam Gereksinimler Listesi}

\begin{lstlisting}[style=kod, language={}, caption={requirements.txt}]
streamlit>=1.35.0
scikit-learn>=1.4.0
pandas>=2.0.0
plotly>=5.18.0
nltk>=3.8.0
numpy>=1.26.0
# Opsiyonel:
# sentence-transformers>=2.7.0
# torch>=2.2.0
\end{lstlisting}

\subsection*{B. Kaynaklar}

\begin{itemize}
    \item Salton, G., \& Buckley, C. (1988). Term-weighting approaches in automatic
          text retrieval. \textit{Information Processing \& Management}, 24(5), 513--523.
    \item Reimers, N., \& Gurevych, I. (2019). Sentence-BERT: Sentence Embeddings
          using Siamese BERT-Networks. \textit{EMNLP 2019}.
    \item Scikit-learn Documentation: \url{https://scikit-learn.org}
    \item Streamlit Documentation: \url{https://docs.streamlit.io}
\end{itemize}

% ─────────────────────────────────────────────────────────────
\end{document}